\section{L7 CPU Architectures}

\subsection{GP-CPU Architectures}

\mult{2}

\begin{definition}{Three Architectures}
    \begin{itemize}
        \item \textbf{Von Neumann}: common I+D memory
        \item \textbf{Harvard}: separate I/D paths
        \item \textbf{Modified Harvard}: Harvard inside the chip, von Neumann outside (IC pins are expensive)
    \end{itemize}
\end{definition}

\begin{definition}{Fundamental Parts}
    \begin{itemize}
        \item cache (speeds transfers into the core)
        \item instruction pipeline + decoder
        \item ALU
    \end{itemize}
\end{definition}

\multend

\begin{KR}{Recipe: Classify a CPU Architecture}
    shared I+D bus $\to$ Von Neumann; fully separate $\to$ Harvard; \emph{split L1 I/D but unified L2/RAM} $\to$ \important{Modified Harvard}. Datasheet cue: separated I/D caches.
\end{KR}

\begin{example2}{Ex.\ (s5) A53 architecture}
    A53: L1-I (2-way) and L1-D (4-way) separate, unified L2 (16-way). What architecture?
     \\ \solution{Solution: }
    \important{Modified Harvard}: Harvard inside (split L1), von Neumann outside (unified L2/memory).
\end{example2}

\subsection{Parallelising a CPU}

\begin{concept}{Four Ways}
    \begin{itemize}
        \item multicore (covered)
        \item multiple FUs (ALUs/multipliers) $\to$ \emph{data} parallelism
        \item multiple pipelines $\to$ \emph{instruction} parallelism
        \item heterogeneous PEs (DSP, GPU)
    \end{itemize}
\end{concept}

\begin{KR}{Recipe: Identify the Parallelism Mechanism}
    multiple FUs (VLIW) $=$ data; multiple pipelines (superscalar) $=$ instruction; idle slots filled by a 2nd thread $=$ SMT; separate cores $=$ multicore/AMP.
\end{KR}

\begin{example2}{Ex.\ Name the four}
    List four ways to add parallelism inside a CPU and the kind each provides.
     \\ \solution{Solution: }
    \important{Multicore}; \important{multiple FUs} (data); \important{multiple pipelines} (instruction); \important{heterogeneous PEs} (DSP/GPU). The ALU is not the only option (e.g.\ FPU/MAC for DSP work).
\end{example2}

\subsection{Multiple FUs: VLIW / DSP}

\mult{2}

\begin{definition}{DSP}
    Application-specific FUs: MAC (.M), data access (.D), GP-ALU (.L), shifter/ALU (.S). Harvard memory.
\end{definition}

\begin{definition}{VLIW}
    Compiler packs serial instructions into one wide word for the FUs.
    \begin{itemize}
        \item needs many registers + bandwidth
        \item large, binary-\important{incompatible} code
    \end{itemize}
\end{definition}

\multend

\begin{concept}{Software Pipelining}
    Overlap loop iterations (prologue / kernel / epilogue) to cut pipeline stalls; with VLIW gives very tight loop kernels.
\end{concept}

\begin{KR}{Recipe: Recognise VLIW / DSP}
    separate program/data bus keeping a MAC fed $\to$ Harvard DSP; one wide instruction the \emph{compiler} packed per FU $\to$ VLIW (+ software pipelining).
\end{KR}

\begin{example2}{Ex.\ (s8) DSP with .M/.D/.L/.S}
    What memory architecture is this? How does execution work?
     \\ \solution{Solution: }
    \important{Harvard} (separate program/data bus). Execution = \important{VLIW}: the compiler packs one op per FU into a wide word, run in parallel (+ software pipelining). Big, incompatible code.
\end{example2}

\subsection{CISC/RISC \& Hazards}

\mult{2}

\begin{definition}{CISC vs RISC}
    \begin{itemize}
        \item \textbf{CISC}: complex instructions, microcoded
        \item \textbf{RISC}: load/store, many simple instructions, pipeline- and compiler-friendly
    \end{itemize}
\end{definition}

\begin{definition}{Hazards}
    \begin{itemize}
        \item structural (resource conflict)
        \item data (dependency)
        \item control (branch)
    \end{itemize}
\end{definition}

\multend

\begin{KR}{Recipe: Spot a Hazard}
    two instructions need the same unit $\to$ structural; one needs another's result $\to$ data; a branch changes the path $\to$ control.
\end{KR}

\begin{example2}{Ex.\ Why RISC pipelines better}
    Why is RISC more pipeline-friendly than CISC? Name the three hazard types.
     \\ \solution{Solution: }
    Uniform load/store instructions fit a fixed pipeline cleanly (CISC needs microcode for variable, complex instructions). Hazards: \important{structural / data / control}.
\end{example2}

\subsection{Superpipeline vs Superscalar}

\mult{2}

\begin{definition}{Superpipelining}
    More / finer pipeline stages $\to$ instructions issued faster.
\end{definition}

\begin{definition}{Superscalar}
    More pipelines $\to$ several instructions per cycle (scheduling becomes the issue).
\end{definition}

\multend

\begin{KR}{Recipe: Tell Them Apart}
    deeper (finer stages, faster issue) $=$ superpipeline; wider (more pipelines, several/cycle) $=$ superscalar.
\end{KR}

\begin{example2}{Ex.\ Distinguish}
    What is the difference between superpipelining and superscalar?
     \\ \solution{Solution: }
    \important{Superpipeline} = more/finer stages (issue faster); \important{superscalar} = more pipelines (several issued per cycle).
\end{example2}

\subsection{In/Out-of-Order Execution}

\begin{definition}{Three Modes}
    \begin{itemize}
        \item in-order issue, in-order execution
        \item in-order issue, out-of-order execution
        \item out-of-order issue, out-of-order execution
    \end{itemize}
\end{definition}

\begin{KR}{Recipe: Identify the Mode}
    Ask whether \emph{issue} and \emph{execution} keep program order. Both ordered = simplest; OoO issue + OoO exec = most aggressive (best unit utilisation, most hardware).
\end{KR}

\begin{example2}{Ex.\ Which is most aggressive?}
    Name the three superscalar scheduling modes; which extracts the most parallelism?
     \\ \solution{Solution: }
    In-order/in-order; in-order issue/OoO exec; \important{OoO issue + OoO exec} (the most aggressive, needs the most scheduling hardware).
\end{example2}

\subsection{Co-processors}

\mult{2}

\begin{definition}{Operation}
    CPU runs recognised instructions; passes others to the co-pro.
    \begin{itemize}
        \item unrecognised $\to$ exception
        \item own clock + memory channels (NEON) $\to$ parallel execution
    \end{itemize}
\end{definition}

\begin{concept}{Interface \& Sync}
    \begin{itemize}
        \item own ISA; instructions via pipeline, data via registers
        \item 8087 + 8086: the 8086 \emph{stalls} on \texttt{FMUL}
        \item future uncertain: costly decoders + compiler upkeep
    \end{itemize}
\end{concept}

\multend

\begin{KR}{Recipe: Reason about a Co-processor}
    CPU offloads unknown instructions; co-pro has its own ISA/clock/memory; the two need synchronisation (CPU may stall); good only for generic ops.
\end{KR}

\begin{example2}{Ex.\ 8087 / 8086}
    How does the CPU handle an instruction it doesn't recognise, and what's the catch with the 8087?
     \\ \solution{Solution: }
    It \important{passes it to the co-pro}; if unrecognised there too $\to$ exception. Catch: the 8086 \important{stalls} while the 8087 performs \texttt{FMUL} (a pipeline/sync effect).
\end{example2}

\subsection{SMT (Hyperthreading)}

\begin{concept}{SMT}
    A core runs 2 threads to use otherwise-idle resources (4 cores $\to$ 8 threads).
    \begin{itemize}
        \item \important{needs a superscalar} core (spare issue slots)
        \item $\sim$30\% speedup (cache-stall dependent); OS sees 2 cores; security issues
    \end{itemize}
\end{concept}

\begin{KR}{Recipe: Reason about SMT}
    SMT fills idle issue slots with a 2nd thread $\Rightarrow$ only worthwhile on a superscalar core; pin a process to a core to control it.
\end{KR}

\begin{example2}{Ex.\ (s19) Why superscalar?}
    Why does SMT require a superscalar core? Typical speedup?
     \\ \solution{Solution: }
    Only a superscalar core has \important{multiple issue slots/pipelines} per cycle to give the 2nd thread; a scalar core has nothing spare. Speedup $\sim$\important{30\%} (more about cache stalls than FU availability).
\end{example2}

\subsection{AMP Scenarios}

\mult{2}

\begin{definition}{Scenarios}
    \begin{itemize}
        \item different ISA (PRU/GPU/DSP)
        \item same ISA (task-driven, master-slave or peer)
        \item same ISA, different clocks / 64-32 (big.LITTLE)
        \item different or no OS
    \end{itemize}
\end{definition}

\begin{concept}{Questions to Answer}
    How do the cores boot? Where is their code? How tightly coupled (sync, data, comms)?
\end{concept}

\multend

\begin{KR}{Recipe: Set Up AMP}
    same ISA $\to$ SMP/big.LITTLE under one OS; \important{different ISA $\to$ separate OS/bare-metal per core + openAMP}. Answer the boot/code/coupling questions first.
\end{KR}

\begin{example2}{Ex.\ (s20) Different ISA under one OS?}
    Can different-ISA cores run under a single OS? Give two AMP scenarios.
     \\ \solution{Solution: }
    \important{No} one kernel is built for one ISA (big.LITTLE works only because it's the \emph{same} ISA). Scenarios: different ISA (specialised PRU/GPU/DSP); same ISA split by task (comms vs I/O); big.LITTLE; different/no OS.
\end{example2}

\subsection{openAMP}

\mult{2}

\begin{definition}{remoteproc}
    Boots a remote core: \texttt{rproc\_boot()} decodes the ELF, places segments in memory, enables resource tables. (Zynq: APU boots Linux via GRUB; RPU R5s start halted.)
\end{definition}

\begin{definition}{rpmsg / virtio}
    IPC: \texttt{rpmsg} channels (source/dest, \texttt{rpmsg\_send} + rx callback) over \texttt{virtio} ring buffers (paravirtualisation split driver).
\end{definition}

\multend

\begin{KR}{Recipe: Couple AMP Cores with openAMP}
    \paragraph{Boot} \texttt{remoteproc} loads + starts the remote ELF.
    \paragraph{Communicate} \texttt{rpmsg} messages over \texttt{virtio} ring buffers.
    \paragraph{Two IPC methods} shared memory + signalling good for asynchronous cores.
\end{KR}

\begin{example2}{Ex.\ Bring up an R5 on the Zynq}
    On the Zynq UltraScale+, how does the APU start an R5 and talk to it?
     \\ \solution{Solution: }
    \important{remoteproc}: \texttt{rproc\_boot} decodes the R5's ELF and boots it (the RPU starts halted). Communication via \important{rpmsg over virtio} ring buffers. Zynq = 4$\times$A53 SMP + 2$\times$R5 lockstep + FPGA accelerators.
\end{example2}

\subsection{Exam FS23 Q2: Architectures (RZ/V2L)}

\begin{concept}{Scenario (appendix)}
    RZ/V2L: Cortex-A55 (1.2\,GHz) + Cortex-M33 (200\,MHz). L1 I/D = 32\,KB each (split), L2 = 0\,KB, L3 = 256\,KB shared; ARMv8.2-A; NEON/FPU; no MMU mentioned; openAMP links A55$\leftrightarrow$M33.
\end{concept}

\begin{KR}{Recipe: Interpret a Cache Hierarchy}
    \paragraph{Datasheet} separate I/D caches $\Rightarrow$ (modified) Harvard.
    \paragraph{Per level} L1 closest (often split I/D); L2 unified (may be 0); L3 shared.
    \paragraph{Missing for design} line size, associativity, total size.
\end{KR}

\begin{example2}{Exam FS23 Q2.1 Architecture (2P)}
    What architecture does this processor represent and why?
     \\ \solution{Solution: }
    \important{Modified Harvard} (1P). Because the data and instruction caches are \emph{separated}, which represents a (modified) Harvard architecture (1P).
\end{example2}

\begin{example2}{Exam FS23 Q2.2 Advantage (2P)}
    What is the advantage of the chosen architecture?
     \\ \solution{Solution: }
    Instruction and data addressing are separated, so fetching one doesn't hinder the other (1P). Different, more fitting cache architectures can be used (1P).
\end{example2}

\begin{example2}{Exam FS23 Q2.3 L1/L2/L3 (max 4P)}
    What do L1, L2 and L3 represent?
     \\ \solution{Solution: }
    Cache architecture (1P). L1 split into data + instruction cache (1P). L2 not existent / 0\,KB (1P). L3 common (shared) cache (1P).
\end{example2}

\begin{example2}{Exam FS23 Q2.5 M33 to A55 (1P)}
    How can the Cortex-M33 communicate with the A55 cores?
     \\ \solution{Solution: }
    Via \important{shared external memory} (1P).
\end{example2}

\begin{example2}{Exam FS23 Q2.6 openAMP IPC (3P)}
    openAMP is in essence two IPC methods which, and why good for M33$\leftrightarrow$A55?
     \\ \solution{Solution: }
    \important{Shared memory} (1P) and \important{signalling} (1P). Because the two cores are \emph{asynchronous} and openAMP suits that asynchronicity (1P).
\end{example2}
